\begin{table}[H]
    \centering
    \caption{Perbandingan Statistik Deskriptif Sampel ($N=1.886$) dan Parameter Nasional Tidak Resmi per Subtes}
    \label{tab:deskriptif}
    \begin{spacing}{1.0}
    \begin{tabular}{lcccccc}
        \toprule
        Subtes & $M$ (Sampel) & $SD$ (Sampel) & $M$ (Nasional) & $SD$ (Nasional) & $\Delta M$ & $\Delta SD$ \\
        \midrule
        PU              & 633,26 & 76,85 & 500,00 & 71,52 & 133,26 & 5,33 \\
        PPU             & 669,96 & 61,50 & 500,00 & 100,00 & 169,96 & -38,50 \\
        PBM             & 680,95 & 76,62 & 500,00 & 100,00 & 180,95 & -23,38 \\
        PK              & 678,53 & 131,90 & 500,00 & 100,00 & 178,53 & 31,90 \\
        LBI Saintek     & 631,63 & 82,31 & 500,00 & 100,00 & 131,63 & -17,69 \\
        LBI Soshum      & 634,69 & 70,11 & 500,00 & 100,00 & 134,69 & -29,89 \\
        LBE             & 680,44 & 86,50 & 500,00 & 100,00 & 180,44 & -13,50 \\
        PM              & 653,23 & 100,01 & 500,00 & 100,00 & 153,23 & 0,01 \\
        \bottomrule
    \end{tabular}
    \end{spacing}
    \vspace{2mm}
    \begin{flushleft}
    \footnotesize\textit{Catatan.} $\Delta M = M_{\text{sampel}} - M_{\text{nasional}}$; $\Delta SD = SD_{\text{sampel}} - SD_{\text{nasional}}$. Parameter nasional tidak resmi bersumber dari tabel internal Panitia Seleksi Nasional Penerimaan Mahasiswa Baru (SNPMB) yang beredar di media sosial dan belum terverifikasi secara resmi (triangulasi parsial pada nilai maksimum PK dan PM; lihat Bagian Metode).
    \end{flushleft}
\end{table}
